% ===== Wrap-up (~4 min, then questions) =====
\partaccent{butter}\notopic

\begin{frame}{Այսօրվա 5 գաղափարը}
  \begin{enumerate}
    \item<1-> Ուղեղը վտանգ է փնտրում \textbf{ամենուր}, ու այսօր «վտանգ» կարող է լինել ամեն ինչ
    \item<2-> Օրինաչափությունների սերն ու պատահականության հանդեպ ատելությունը ծնում են սնահավատություններ ու դավադրության տեսություններ (\textbf{ապոֆենիա})
    \item<3-> Ֆոբիաները սովորած վախեր են, որոնք \textbf{ինքն իրենց ապացուցում} են. խուսափելը դրանք պահպանում է
    \item<4-> Վախը կարող է հաճույք լինել. երբ վտանգն անցնում է, պարգևի ուղին ուժեղ «պարգև» է տալիս, եթե \textbf{մի քիչ} վերահսկողություն ունես
    \item<5-> Մեկ քննադատությունն ավելի ուժեղ է, քան շատ գովեստներ. կորտիզոլ, նորություն, եսասեր ուղեղ, կասկածելի քաղաքավարություն
  \end{enumerate}
\end{frame}

\begin{frame}{Փորձիր տանը}
  \begin{columns}[T]
    \begin{column}{0.31\textwidth}
      \begin{card}[colback=lavender!35!cardwhite]
        \centering\bigemoji[1cm]{clover}\par\smallskip
        \raggedright Մեկ շաբաթ նկատիր քո փոքրիկ «ծեսերը». թակե՞լ ես փայտին, շրջանցե՞լ ես սև կատվին։ Երբ ու ինչու
      \end{card}
    \end{column}
    \begin{column}{0.31\textwidth}
      \begin{card}[colback=peach!35!cardwhite]
        \centering\bigemoji[1cm]{film}\par\smallskip
        \raggedright Սարսափ ֆիլմ նայելիս նկատիր՝ երբ է ավելի վախենալի. երբ հրեշը երևում է, թե՞ երբ դեռ չի երևում
      \end{card}
    \end{column}
    \begin{column}{0.31\textwidth}
      \begin{card}[colback=brainpink!25!cardwhite]
        \centering\bigemoji[1cm]{memo}\par\smallskip
        \raggedright Մի օր գրի առ քեզ ասված գովեստներն ու քննադատությունները։ Մյուս օրը ո՞րն ես ավելի լավ հիշում
      \end{card}
    \end{column}
  \end{columns}
\end{frame}

\begin{frame}{Հաջորդ վեբինարը}
  \begin{columns}[c]
    \begin{column}{0.62\textwidth}
      \chip[lavender]{graduation}{4. Կարծում ես՝ խելացի՞ ես}\par\medskip
      \begin{itemize}
        \item «Իմ IQ-ն 270 է». ինչու է խելքը չափելն այդքան դժվար
        \item «Որտե՞ղ է տաբատդ, պրոֆեսոր». ինչու են խելացիները հիմարություններ անում
        \item Ինչու են խելացի մարդիկ հաճախ վեճում պարտվում
        \item Խաչբառերն ուղեղդ չեն սրում
        \item Բարձրահասակներն ավելի խելացի՞ են, ու որքանով է խելքը ժառանգվում
      \end{itemize}
    \end{column}
    \begin{column}{0.32\textwidth}
      \centering\mascot[2.8cm]{thinking}
    \end{column}
  \end{columns}
\end{frame}

\begin{frame}[plain]
  \begin{tikzpicture}[remember picture, overlay]
    \fill[lavender!60!cream] (current page.south west) rectangle (current page.north east);
    \fill[pattern={Dots[distance=4.4mm, radius=0.28mm]}, pattern color=plum!10!lavender!60!cream]
      (current page.south west) rectangle (current page.north east);
    \node at ([xshift=-3.2cm]current page.center) {\mascot[4.4cm]{happy}};
    \node[anchor=west, text=plum] at ([xshift=0.2cm, yshift=0.6cm]current page.center)
      {\titlefont\fontsize{40}{42}\selectfont Հարցեր};
    \node[anchor=west, text=plum, font=\mediumfont\large, align=left] at ([xshift=0.2cm, yshift=-1.0cm]current page.center)
      {Գրեք չաթում կամ\\ միացրեք խոսափողը \emoji{megaphone}};
    \node[anchor=west, text=plum!70!cream, font=\small] at ([xshift=0.2cm, yshift=-2.1cm]current page.center)
      {Դին Բըրնեթ, «The Idiot Brain», 2016};
  \end{tikzpicture}
\end{frame}
